\documentclass[11pt]{article}
\usepackage[T1]{fontenc}
\usepackage{amsmath,amssymb,amsthm}
\usepackage[margin=1in]{geometry}
\usepackage{graphicx}
\usepackage{microtype}
\usepackage{xcolor}
\usepackage[hidelinks]{hyperref}

\newtheorem{theorem}{Theorem}
\newtheorem{proposition}[theorem]{Proposition}
\newcommand{\vol}{\operatorname{vol}}
\newcommand{\B}{B_2}

\title{A determinant-one counterexample to the proposed\\
strong-\((B)\) equality classification}
\author{Automated Math Discovery Run \texttt{fa\_banach\_001}, lane 8}
\date{11 August 2026}

\begin{document}
\maketitle

\begin{center}
\fcolorbox{black}{red!5}{\parbox{0.88\linewidth}{
\textbf{Result.} The equality characterization in Conjecture~4.1 of
Artstein-Avidan--Katzin is false.  There is a centrally symmetric
unconditional convex body \(K\subset\mathbb R^3\) and a nonzero,
trace-zero, nonscalar diagonal matrix \(\Lambda\) for which
\(\vol(e^{s\Lambda}K\cap\B^3)\) is constant on a nontrivial interval,
although neither \(K\subset\B^3\) nor \(\B^3\subset K\).
}}
\end{center}

\section{Source conjecture}

Conjecture~4.1 of \cite{AAK} asserts that, for a convex body
\(K\subset\mathbb R^n\) and diagonal \(\Lambda\),
\[
 \phi(t)=\vol_n(e^{t\Lambda}K\cap\B^n)
\]
is log-concave.  In particular it proposes
\begin{equation}\label{eq:midpoint}
 \vol_n(e^{t\Lambda/2}K\cap\B^n)^2
 \ge
 \vol_n(e^{t\Lambda}K\cap\B^n)\,
 \vol_n(K\cap\B^n).
\end{equation}
It further states that equality occurs if and only if
\(K\subset\B^n\), or \(\B^n\subset K\), or
\(\Lambda=\lambda I_n\).  The paper uses this rigidity clause, with
\(\operatorname{tr}\Lambda=0\), to deduce uniqueness of a maximum
intersection ellipsoid in the intermediate-radius regime.

\begin{figure}[ht]
\centering
\includegraphics[width=0.96\linewidth,height=0.42\textheight,keepaspectratio]{figures/open_problem_crop.png}
\caption{Conjecture~4.1 and Proposition~4.2 in arXiv:1612.01128,
PDF pages 8--9.}
\end{figure}

\section{Idea of the counterexample}

Let one coordinate of a box cut the unit ball and make the other coordinate
widths larger than the ball.  A diagonal deformation supported only on the
oversized coordinates is invisible to intersection with the ball for as
long as those widths remain at least one.  Two such redundant coordinates
can expand and contract reciprocally, making the deformation nontrivial and
trace zero while leaving the intersection exactly fixed.

\section{The counterexample}

\begin{theorem}\label{thm:counterexample}
Set
\[
 K=[-\tfrac12,\tfrac12]\times[-2,2]\times[-3,3]
 \quad\text{and}\quad
 \Lambda=\operatorname{diag}(0,1,-1).
\]
Then, for every \(s\in[0,1]\),
\begin{equation}\label{eq:constant-set}
 e^{s\Lambda}K\cap\B^3
 =\{x\in\B^3:|x_1|\le\tfrac12\}.
\end{equation}
Consequently \eqref{eq:midpoint} is an equality at \(t=1\), although none
of its three proposed equality conditions holds.  Moreover
\(\operatorname{tr}\Lambda=0\).
\end{theorem}

\begin{proof}
For \(s\in[0,1]\),
\[
 e^{s\Lambda}K
 =[-\tfrac12,\tfrac12]\times[-2e^s,2e^s]
   \times[-3e^{-s},3e^{-s}].
\]
The second half-width satisfies \(2e^s\ge2>1\), and the third satisfies
\(3e^{-s}\ge3/e>1\).  Since every \(x\in\B^3\) has
\(|x_2|,|x_3|\le1\), those two box constraints are redundant inside the
ball.  This proves \eqref{eq:constant-set}.

Writing \(a=1/2\) and integrating the disk cross-sections gives the common
volume
\[
 \phi(s)=\int_{-a}^{a}\pi(1-u^2)\,du
 =2\pi\left(a-\frac{a^3}{3}\right)=\frac{11\pi}{12}.
\]
Thus \(\phi(1/2)^2=\phi(1)\phi(0)\).

The point \((0,2,0)\) belongs to \(K\setminus\B^3\), while
\((1,0,0)\) belongs to \(\B^3\setminus K\).  Finally, \(\Lambda\) has
three unequal diagonal entries, so it is not scalar, while
\(\operatorname{tr}\Lambda=0+1-1=0\).
\end{proof}

\section{A robust family of equality cases}

\begin{proposition}\label{prop:family}
Let \(0<a<1\), let \(I\subset\mathbb R\) be an interval, and put
\[
 K=[-a,a]\times\prod_{j=2}^n[-b_j,b_j],\qquad
 \Lambda=\operatorname{diag}(0,\lambda_2,\ldots,\lambda_n).
\]
If \(b_j e^{s\lambda_j}\ge1\) for every \(s\in I\) and every
\(j\ge2\), then
\[
 e^{s\Lambda}K\cap\B^n
 =\{x\in\B^n:|x_1|\le a\}
 \qquad(s\in I).
\]
In particular, in every dimension \(n\ge3\) one may choose a nonzero
trace-zero nonscalar \(\Lambda\) and obtain equality on any prescribed
compact interval after taking the \(b_j\)'s sufficiently large.
\end{proposition}

\begin{proof}
Inside \(\B^n\), every coordinate has absolute value at most one.  The
hypothesis therefore makes all constraints except \(|x_1|\le a\)
redundant, proving the set identity.  For the final assertion, take two
indices \(j,k\ge2\), set \(\lambda_j=c\), \(\lambda_k=-c\) with
\(c\ne0\), and set the remaining eigenvalues to zero.  On a compact
interval, finite \(b_j,b_k\) can dominate the largest reciprocal scaling.
\end{proof}

This family shows that the failure is structural.  It is not caused by
choosing \(t=0\), by a scalar dilation, by non-central symmetry, or by a
volume-changing deformation.  The intersection set itself, not merely its
volume, remains fixed on a whole interval.

\section{Why this does not yet settle ellipsoid uniqueness}

Because \(\operatorname{tr}\Lambda=0\), change of variables gives
\[
 \vol(e^{s\Lambda}K\cap\B^3)
 =\vol(K\cap e^{-s\Lambda}\B^3).
\]
Thus a continuum of distinct equal-volume ellipsoids has the same overlap
with \(K\).  To refute uniqueness of the \emph{maximizing} ellipsoid, one
would still have to prove that this common value is the global maximum.  It
is not so for the displayed example.

Indeed, use the source's first-variation formula at the unit ball with
\(A=\operatorname{diag}(-2,1,1)\).  On the sphere,
\(S^2\cap K=\{x\in S^2:|x_1|\le a\}\), \(a=1/2\), because the other box
constraints are redundant.  Since spherical area disintegrates as
\(dS=2\pi\,du\) with \(u=x_1\),
\begin{align*}
 \left.\frac d{dt}\right|_{t=0}
 \vol(K\cap e^{tA}\B^3)
 &=\int_{S^2\cap K}\langle x,Ax\rangle\,dS\\
 &=4\pi a-3\left(\frac{4\pi a^3}{3}\right)
 =4\pi a(1-a^2)>0.
\end{align*}
So the unit ball is not even a local maximizer.  This closes the most direct
upgrade route without overstating the result: the equality rigidity in
Conjecture~4.1 is false, while the log-concavity inequality and the paper's
intermediate-radius uniqueness question remain open here.

\section{Verification and novelty notes}

The proof is an exact coordinate calculation.  The accompanying script
checks all endpoint widths on a 1001-point grid, the exact cross-sectional
volume, and the positive first variation; the grid is only a transcription
check, since monotonicity proves the bounds analytically.

A bounded official-arXiv search through 11 August 2026 used arXiv id
1612.01128 and combinations of ``strong B equality'', ``uniform measure
ball equality'', and ``maximal intersection ellipsoid uniqueness''.  It
found the source and related strong-\(B\) literature but no later answer to
this exact equality claim.  This is not an exhaustive literature review.

\begin{thebibliography}{9}
\bibitem{AAK}
S.~Artstein-Avidan and D.~Katzin,
\emph{Isotropic Measures and Maximizing Ellipsoids: Between John and
Loewner}, arXiv:1612.01128.
\end{thebibliography}

\end{document}
